\chapter{Introduction}
In this rapport we will setup and use the Michelson Interferometer,
to measure the expansion of a piezo electric material, when subject to varying voltages.
\section{The Michelson Interferometer}
The Michelson Interferometer is a device, in which light is led through a beam splitter, 
afterwhich it moves in two paths of different lengths, before being reflected back towards the beam splitter by a mirror.
When the light hits the beam splitter again, they are reflected and transmitted such that they move in the save direction towards a sensor.
The light that is then measured depends of the difference in lengths that the light moved, because the light hits the sensor with different phases, 
and will therefore interfere, as the name interferometer suggests.
\section{Interference}
The light from the laser has a high degree of spatial coherence because of the resonance in the laser cavity. 
This means that the phase is constant across the beams wavefront. 
When the laser beam is split by the beamsplitter and reflected back to meet again the phases should lineup to interfere. 
The difference in distance of the two mirrors change the interference from destructive to constructive and in between. 
